\documentclass[11pt]{article}

% =====================================================================
%  Dynamic Lagrangian Frequency--Time (DLFT) in admittance / FBS form
%  Companion derivation for the pyhbm implementation (class FBS_DLFT).
%
%  Notation follows Krack & Gross, "Harmonic Balance for Nonlinear
%  Vibration Problems" (2019); where it diverges we follow
%  Vadcard, Batailly & Thouverez, J. Sound Vib. 531 (2022) 116950.
%  FBS assembly follows van der Seijs, PhD thesis (2016), Ch. 3
%  (LM-FBS, Eqs. 3.20-3.26) and Soleimani & Aage, CMAME 439 (2025)
%  117882 (closed-form nonlinear FBS, Eq. 24).
% =====================================================================

\usepackage[a4paper,margin=2.4cm]{geometry}
\usepackage{amsmath,amssymb,amsthm}
\usepackage{bm}
\usepackage{booktabs}
\usepackage{enumitem}
\usepackage[colorlinks=true,linkcolor=blue,citecolor=blue]{hyperref}

% ---- macros -------------------------------------------------------
\newcommand{\R}{\mathbb{R}}
\newcommand{\C}{\mathbb{C}}
\newcommand{\T}{^{\mathsf{T}}}
\newcommand{\inv}{^{-1}}
\newcommand{\Mb}{\bm{M}}
\newcommand{\Cb}{\bm{C}}
\newcommand{\Kb}{\bm{K}}
\newcommand{\Zb}{\bm{Z}}
\newcommand{\Yb}{\bm{Y}}
\newcommand{\Bb}{\bm{B}}
\newcommand{\Ib}{\bm{I}}
\newcommand{\ub}{\bm{u}}
\newcommand{\fb}{\bm{f}}
\newcommand{\xb}{\bm{x}}
\newcommand{\gb}{\bm{g}}
\newcommand{\lamb}{\bm{\lambda}}
\newcommand{\Gam}{\bm{\Gamma}}
\newcommand{\diag}{\operatorname{diag}}
\newcommand{\blkdiag}{\operatorname{blkdiag}}
\newcommand{\fadm}{\bm{f}_{\!\mathrm{adm}}}
\newcommand{\xr}{\tilde{\bm{x}}_{\mathrm r}}      % multiharmonic relative displacement
\newcommand{\lamt}{\tilde{\bm{\lambda}}}          % multiharmonic contact force
\newcommand{\Yr}{\bm{Y}_{\mathrm r}}
\newcommand{\Zr}{\bm{Z}_{\mathrm r}}
\newcommand{\Real}{\operatorname{Re}}
\newcommand{\Imag}{\operatorname{Im}}

\title{\bf Dynamic Lagrangian Frequency--Time (DLFT) in admittance
notation for Frequency-Based Substructuring}
\author{Companion note for \texttt{pyhbm} -- class \texttt{FBS\_DLFT}}
\date{\today}

\begin{document}
\maketitle

\begin{abstract}
This note derives the Dynamic Lagrangian Frequency--Time (DLFT)
harmonic-balance method of Vadcard \emph{et al.} (2022) directly in
\emph{admittance} (receptance/FRF) notation and embeds it in a
Lagrange-multiplier Frequency-Based Substructuring (LM-FBS) assembly
that uses signed Boolean matrices $\Bb_{\mathrm c},\Bb_{\mathrm e}$ on a
block-diagonal substructure admittance $\Yb=\blkdiag(\Yb^1,\Yb^2,\dots)$.
The closed-form coupling of Soleimani \& Aage (2025, Eq.~24) is recovered
as the two-substructure special case. We give the residual, the
\emph{full} analytic Jacobian $\partial\bm r/\partial\xr$ and
$\partial\bm r/\partial\omega$ (including the conditional contact term and
the $\partial\Yb/\partial\omega$ chain rule), state exactly where the
DFT/IDFT enter, and answer the algorithmic question: \emph{the primary
Newton unknown is the multiharmonic vector of relative-displacement
Fourier coefficients $\xr$, not the contact forces $\lamt$.} The
derivation is the specification for the \texttt{FBS\_DLFT} class and the
predictor--corrector continuation loop in \texttt{pyhbm}.
\end{abstract}

\tableofcontents

% =====================================================================
\section{Notation and harmonic-balance preliminaries}
% =====================================================================
Consider a linear elastic substructure with mass, damping and stiffness
matrices $\Mb,\Cb,\Kb\in\R^{d\times d}$ driven at the fundamental angular
frequency $\omega$. With the (real) Fourier ansatz for a $T=2\pi/\omega$
periodic displacement,
\begin{equation}
\ub(t)\;=\;\Real\!\sum_{n\in\mathcal H}\hat\ub_n\,e^{\,\mathrm i n\omega t},
\qquad \mathcal H=\{0,1,\dots,N_H\},\quad N_h=|\mathcal H|,
\label{eq:ansatz}
\end{equation}
each harmonic $n$ obeys the algebraic balance
\begin{equation}
\Zb_n(\omega)\,\hat\ub_n \;=\; \hat\fb_n,
\qquad
\Zb_n(\omega)\;=\;\Kb-(n\omega)^2\Mb+\mathrm i\,n\omega\,\Cb \;\in\;\C^{d\times d},
\label{eq:Zn}
\end{equation}
where $\Zb_n$ is the \emph{dynamic stiffness} (impedance) and
$\Yb_n(\omega)=\Zb_n(\omega)\inv$ the \emph{admittance} (receptance/FRF) of
harmonic $n$. Stacking harmonics,
$\xr,\hat\fb\in\C^{N_h d}$, the block-diagonal multiharmonic operators are
\begin{equation}
\Zb(\omega)=\blkdiag_{n\in\mathcal H}\Zb_n(\omega),
\qquad
\Yb(\omega)=\Zb(\omega)\inv=\blkdiag_{n\in\mathcal H}\Yb_n(\omega).
\label{eq:ZYblock}
\end{equation}
For a nonlinear system with internal force $\fb_{\mathrm{nl}}$ the
multiharmonic balance residual reads, in impedance form
(Vadcard Eq.~23; Krack \S2),
\begin{equation}
\mathcal H(\hat\ub,\omega)\;=\;\Zb(\omega)\,\hat\ub
+\tilde\fb_{\mathrm{nl}}(\hat\ub)-\tilde\fb_{\mathrm{ext}}\;=\;\bm 0.
\label{eq:HBimp}
\end{equation}

\paragraph{DFT/IDFT operators.}
Nonlinear and contact forces are evaluated in the time domain on
$N_t$ equispaced samples $\tau_i=2\pi i/N_t$, $i=0,\dots,N_t-1$
($N_t$ over-samples the band $\mathcal H$). Let
\begin{equation}
\Gam:\ \C^{N_h d}\!\to\R^{N_t d}\ \text{(synthesis / IDFT)},
\qquad
\Gam^{+}:\ \R^{N_t d}\!\to\C^{N_h d}\ \text{(analysis / DFT)},
\label{eq:dft}
\end{equation}
so that a band-limited real signal $\bm s(t_i)=\Gam\,\hat{\bm s}$ and
$\hat{\bm s}=\Gam^{+}\bm s$, with $\Gam^{+}\Gam=\Ib$ on the retained band.
Acting on the real/imaginary split $\hat{\bm s}=\hat{\bm s}_R+\mathrm i\hat{\bm s}_I$,
the (real) entries of $\Gam,\Gam^{+}$ for one DOF are
\begin{equation}
\Gam_{i,n}^{(R)}=\tfrac{c_n}{N_t}\cos(n\tau_i),\quad
\Gam_{i,n}^{(I)}=-\tfrac{c_n}{N_t}\sin(n\tau_i),\quad
\Gam^{+(R)}_{n,i}=\cos(n\tau_i),\quad
\Gam^{+(I)}_{n,i}=-\sin(n\tau_i),
\label{eq:dftentries}
\end{equation}
with $c_0=1$, $c_{n>0}=2$ (the \texttt{numpy.rfft}/\texttt{irfft}
convention used in \texttt{pyhbm}). The distinction
``$\Gam$ acts on the \emph{real} pair $(\hat{\bm s}_R,\hat{\bm s}_I)$''
is essential for the Jacobian (\S\ref{sec:jac}).

% =====================================================================
\section{LM-FBS assembly and condensation to the relative interface DOF}
% =====================================================================
\subsection{Dual (Lagrange-multiplier) assembly}
Let substructures $s=1,\dots,N_s$ have admittances $\Yb^s(\omega)$.
Concatenate the global, \emph{uncoupled} quantities
\begin{equation}
\ub=\operatorname{col}(\ub^1,\dots,\ub^{N_s}),\quad
\fb=\operatorname{col}(\fb^1,\dots,\fb^{N_s}),\quad
\Yb=\blkdiag(\Yb^1,\dots,\Yb^{N_s}),
\label{eq:concat}
\end{equation}
(each understood per harmonic, or block-diagonal over $\mathcal H$ as in
\eqref{eq:ZYblock}). Following van der Seijs (Eqs.~3.20, 3.25), coupling is
imposed through a \emph{compatibility} Boolean $\Bb_{\mathrm c}$ and an
\emph{equilibrium} Boolean $\Bb_{\mathrm e}$: the connecting forces are
$\bm g=-\Bb_{\mathrm e}\T\lamb$ and compatibility reads
$\Bb_{\mathrm c}\ub=\bm 0$. The dual system is
\begin{equation}
\ub=\Yb\bigl(\fb-\Bb_{\mathrm e}\T\lamb\bigr),
\qquad \Bb_{\mathrm c}\ub=\bm 0 .
\label{eq:dual}
\end{equation}
For \emph{collocated, conforming} interfaces (the case implemented here)
$\Bb_{\mathrm c}=\Bb_{\mathrm e}=:\Bb$ is the signed Boolean with $\pm1$ on
each connected DOF pair. Eliminating $\lamb$ for a \emph{rigid} joint gives
the classical LM-FBS one-liner (van der Seijs Eqs.~3.21--3.23)
\begin{equation}
\lamb=\bigl(\Bb_{\mathrm c}\Yb\Bb_{\mathrm e}\T\bigr)\inv \Bb_{\mathrm c}\Yb\fb,
\qquad
\tilde\Yb=\Bigl[\Ib-\Yb\Bb_{\mathrm e}\T(\Bb_{\mathrm c}\Yb\Bb_{\mathrm e}\T)\inv\Bb_{\mathrm c}\Bigr]\Yb .
\label{eq:lmfbs}
\end{equation}

\subsection{Nonlinear interface: condensation onto the relative DOF}
\label{sec:cond}
With a \emph{nonlinear} interface the rigid constraint
$\Bb_{\mathrm c}\ub=\bm0$ is replaced by an interface constitutive law: the
multiplier $\lamt$ becomes the (unknown) contact force that depends on the
relative interface motion. Define the multiharmonic
\emph{relative interface displacement}
\begin{equation}
\xr \;:=\; \Bb_{\mathrm c}\,\hat\ub \;\in\;\C^{N_h\,n_{\mathrm{int}}},
\label{eq:xrel}
\end{equation}
($n_{\mathrm{int}}$ interface DOF pairs). Substituting the first line of
\eqref{eq:dual} into \eqref{eq:xrel},
\begin{equation}
\xr=\Bb_{\mathrm c}\Yb\fb-\Bb_{\mathrm c}\Yb\Bb_{\mathrm e}\T\lamt
=\fadm-\Yr\,\lamt,
\label{eq:cond}
\end{equation}
where we name the \emph{interface admittance} and the \emph{admittance
(free) excitation response}
\begin{equation}
\boxed{\;\Yr:=\Bb_{\mathrm c}\Yb\Bb_{\mathrm e}\T\;}\quad(=\Yb^1+\Yb^2\ \text{collocated, 2 subs.}),
\qquad
\boxed{\;\fadm:=\Bb_{\mathrm c}\Yb\,\tilde\fb_{\mathrm{ext}}\;}.
\label{eq:YrFadm}
\end{equation}
Equation \eqref{eq:cond} rearranges to the \textbf{admittance FBS residual}
\begin{equation}
\boxed{\;\bm r(\xr)\;=\;\xr+\Yr\,\lamt-\fadm\;=\;\bm 0\;}
\label{eq:residual}
\end{equation}
which is the master equation of this note. It is the admittance
counterpart of Vadcard's impedance condensation $\Zr\xr+\lamt-\tilde\fb_{\mathrm{ex,r}}=\bm0$
(his Eq.~37): multiply \eqref{eq:residual} from the left by
$\Zr:=\Yr\inv$ and use $\Zr\fadm=\tilde\fb_{\mathrm{ex,r}}$ to recover it.
Hence the impedance and admittance residuals share the same root; they
differ only by the left factor $\Yr$, i.e.\ by conditioning
(see \S\ref{sec:cond-number}).

\subsection{Recovery of Soleimani's closed form (two substructures)}
For $N_s=2$ collocated substructures, $\Bb_{\mathrm c}=\Bb_{\mathrm e}=[\,\Ib\ \;-\Ib\,]$ on
the boundary DOFs, $\Yb=\blkdiag(\Yb^1,\Yb^2)$, so
\begin{equation}
\Yr=\Bb\Yb\Bb\T=\Yb^1_{bb}+\Yb^2_{bb},\qquad
\fadm=\Yb^1_{be}\hat\fb^1_e-\Yb^2_{be}\hat\fb^2_e,
\label{eq:soltwo}
\end{equation}
and \eqref{eq:residual} reads
$\xr+(\Yb^1_{bb}+\Yb^2_{bb})\lamt-(\Yb^1_{be}\hat\fb^1_e-\Yb^2_{be}\hat\fb^2_e)=\bm0$,
i.e.\ exactly Soleimani \& Aage Eq.~(24) with $\lamt=\bm f_N$ (their
receptance $\bm H\equiv\Yb$). The grounded/unilateral case
(``one substructure rigid'') sets $\Yb^2\to\bm0$, giving
$\Yr=\Yb^1_{bb}$, $\fadm=\Yb^1_{be}\hat\fb^1_e$ -- the wall test case in
\texttt{examples/fbs\_dnft\_unilateral\_spring}.

% =====================================================================
\section{Unilateral contact and the DLFT slaving of $\lamt$}
% =====================================================================
\subsection{Constraints}
The interface is a unilateral (compression-only) contact. With the gap
(clearance) function and contact pressure $\lambda(t)$ the
Hertz--Signorini--Moreau / unilateral contact conditions (UCC) are
\begin{equation}
g(t)\ge 0,\qquad \lambda(t)\ge 0,\qquad g(t)\,\lambda(t)=0 .
\label{eq:ucc}
\end{equation}
We use the (implemented) convention of a wall at relative offset $g_0$:
the scalar relative coordinate $x_{\mathrm r}(t)=(\Gam\xr)(t)$ may not pass
the wall, so
\begin{equation}
g(t)=g_0-x_{\mathrm r}(t),\qquad
\text{penetration } p(t)=x_{\mathrm r}(t)-g_0=-g(t).
\label{eq:gap}
\end{equation}
(Vadcard writes $g=x_{\mathrm r}+g_0$ with the opposite sign of
$x_{\mathrm r}$; \eqref{eq:gap} is the mirror convention and is physically
equivalent for a wall at $+g_0$.)

\subsection{Prediction step (admittance form)}
\label{sec:pred}
The contact force is not an independent unknown; it is \emph{slaved} to
$\xr$ by a prediction--correction (this is what makes DLFT cheap -- it does
\emph{not} double the unknowns with explicit multipliers). Isolating the
multiplier in the equilibrium relation and adding a gap penalty with
factor $\varepsilon>0$ (Vadcard Eq.~40), the prediction is
\begin{equation}
\lamt_{\mathrm p}(\xr)
=\Zr\bigl(\fadm-\xr\bigr)\;-\;\varepsilon\,\tilde{\bm g}
=\underbrace{\Zr\bigl(\fadm-\xr\bigr)}_{\text{balancing force}}
+\;\varepsilon\,(\xr-\tilde{\bm g}_0),
\qquad \Zr=\Yr\inv .
\label{eq:pred}
\end{equation}
The first term is the force that exactly satisfies the balance
\eqref{eq:residual} at the current $\xr$ (note
$\Zr(\fadm-\xr)=\tilde\fb_{\mathrm{ex,r}}-\Zr\xr$, recovering Vadcard's
$\tilde\fb_{\mathrm{ex,r}}-\Zr\xr$). Forming $\Zr=\Yr\inv$ from the
substructure admittances is the one place the admittance assembly must
\emph{invert} the interface admittance $\Yr=\Bb_{\mathrm c}\Yb\Bb_{\mathrm e}\T$ -- the same
matrix whose conditioning van der Seijs (after Eq.~3.26) calls ``highly
defining for the success of the coupling.''

\paragraph{Where the IDFT enters.} The UCC \eqref{eq:ucc} are pointwise in
time, so the prediction is taken to the time domain,
\begin{equation}
\lambda_{\mathrm p}(t_i)=\bigl(\Gam\,\lamt_{\mathrm p}\bigr)_i
=\bigl(\Gam\,\Zr(\fadm-\xr)\bigr)_i+\varepsilon\bigl(x_{\mathrm r}(t_i)-g_0\bigr).
\label{eq:predtime}
\end{equation}
Interpretation (Vadcard \S2.4.2): where the gap is \emph{violated}
($p>0$) the penalty drives $\lambda_{\mathrm p}$ strongly positive
(contact detected); where the bodies \emph{separate} ($p<0$, i.e.\
$g>0$) it drives $\lambda_{\mathrm p}$ negative (to be removed below);
in persistent contact ($g\approx0$) the penalty is negligible.

\subsection{Correction step (UCC enforcement) and where the DFT enters}
\label{sec:corr}
The predicted pressure violates $\lambda\ge0$ during separation. The
correction enforces positivity \emph{pointwise in time} (Vadcard Eq.~42),
\begin{equation}
\lambda(t_i)=\max\!\bigl(0,\ \lambda_{\mathrm p}(t_i)\bigr)
=\begin{cases}\lambda_{\mathrm p}(t_i)& \lambda_{\mathrm p}(t_i)>0\ \text{(contact)}\\[2pt]
0 & \lambda_{\mathrm p}(t_i)\le 0\ \text{(separation)}\end{cases}
\label{eq:corr}
\end{equation}
and is transformed back to the frequency domain by the \emph{DFT}:
\begin{equation}
\lamt(\xr)=\Gam^{+}\bigl\{\max(0,\lambda_{\mathrm p}(t_i))\bigr\}_{i=1\dots N_t}.
\label{eq:lamtilde}
\end{equation}
Equations \eqref{eq:pred}--\eqref{eq:lamtilde} define the map
$\xr\mapsto\lamt(\xr)$; substituting into \eqref{eq:residual} gives the
\emph{single}-argument residual $\bm r(\xr)$ that the corrector solves.
The contact set is encoded by the time-domain Boolean mask
\begin{equation}
m_i=\mathbf 1\!\left[\lambda_{\mathrm p}(t_i)>0\right]\in\{0,1\},
\qquad \bm m=\diag(m_1,\dots,m_{N_t}) .
\label{eq:mask}
\end{equation}

% =====================================================================
\section{The primary Newton unknown}
% =====================================================================
\noindent\textbf{Claim.} \emph{The algorithmic Newton unknown is the
multiharmonic relative displacement $\xr$ (its $2N_h n_{\mathrm{int}}$ real
DOF), not the contact force $\lamt$.}

\smallskip\noindent\emph{Justification from the residual/Jacobian structure
(not convention):}
\begin{enumerate}[leftmargin=1.6em,itemsep=2pt]
\item After \eqref{eq:pred}--\eqref{eq:lamtilde}, $\lamt$ is an explicit
\emph{function} $\lamt(\xr)$. The residual \eqref{eq:residual} therefore
depends on $\xr$ \emph{only}; the root-finding lives in
$\R^{2N_h n_{\mathrm{int}}}$ and the Jacobian
$\partial\bm r/\partial\xr$ (\S\ref{sec:jac}) is \emph{square} in that
space.
\item Promoting $\lamt$ to an independent unknown would require appending
the prediction equation and the complementarity \eqref{eq:ucc} as extra
equations, \emph{doubling} the system and re-introducing the explicit
Lagrange multipliers DLFT is designed to avoid (Vadcard \S2.4.1; this is
the very difference between DLFT and a wet/saddle-point Lagrangian, cf.\
van der Seijs Eq.~3.28). The admittance form makes this concrete: $\Yr$
in \eqref{eq:residual} already \emph{is} the condensed interface operator,
so no auxiliary force DOF are needed.
\item Consequently the continuation tangent and step are taken in
$(\xr,\omega)$. The contact force is never a state: it is recomputed from
$\xr$ at \emph{every} residual evaluation.
\end{enumerate}

\noindent\textbf{Distinguish two notions.}
\emph{(i) Choice of unknowns (algorithmic):} $\xr$, as argued above.
\emph{(ii) Initial guess at each continuation step (predictor):} the
tangent predictor extrapolates the converged $(\xr,\omega)$ of the
previous point; equivalently the linear relative response
$\xr^{(0)}=\fadm$ (no contact) is a natural cold start. \textbf{In both
cases the carried quantity is the relative-displacement Fourier
coefficients, never the contact forces} -- $\lamt$ is regenerated by
\eqref{eq:pred}--\eqref{eq:lamtilde}.

% =====================================================================
\section{Analytic Jacobian}
\label{sec:jac}
% =====================================================================
The corrector and the arc-length predictor need
$\partial\bm r/\partial\xr$ and $\partial\bm r/\partial\omega$. Both follow
from \eqref{eq:residual} by the chain rule through the slaving map.

\subsection{Derivative w.r.t.\ the relative-displacement coefficients}
From \eqref{eq:residual},
\begin{equation}
\frac{\partial\bm r}{\partial\xr}
=\Ib+\Yr\,\frac{\partial\lamt}{\partial\xr}.
\label{eq:drdx}
\end{equation}
The slaving map \eqref{eq:pred}--\eqref{eq:lamtilde} differentiates in
three stages.

\paragraph{(a) Prediction (frequency, holomorphic).} From \eqref{eq:pred},
$\lamt_{\mathrm p}=(\Zr\fadm-\varepsilon\tilde{\bm g}_0)+(\varepsilon\Ib-\Zr)\xr$, hence
\begin{equation}
\frac{\partial\lamt_{\mathrm p}}{\partial\xr}=\varepsilon\Ib-\Zr
\qquad(\text{Vadcard A.3: }\ -\Zr-\varepsilon\Ib\ \text{in his sign convention}).
\label{eq:dpred}
\end{equation}
This block is \emph{complex-linear} (holomorphic): in the real/imaginary
(RI) split it has the structure $\bigl[\begin{smallmatrix}\Real&-\Imag\\\Imag&\Real\end{smallmatrix}\bigr]$.

\paragraph{(b) Correction (time, conditional / non-holomorphic).} The
$\max(0,\cdot)$ in \eqref{eq:corr} has derivative $m_i$ pointwise
(Vadcard A.4): keep the predicted sensitivity where there is contact, zero
it where there is separation. With the mask \eqref{eq:mask},
\begin{equation}
\frac{\partial\lambda(t_i)}{\partial\xr}
=\begin{cases}\dfrac{\partial\lambda_{\mathrm p}(t_i)}{\partial\xr}&\lambda_{\mathrm p}(t_i)>0\\[6pt]
\bm 0&\lambda_{\mathrm p}(t_i)\le0\end{cases}
\;=\;\bm m\,\Gam\,\frac{\partial\lamt_{\mathrm p}}{\partial\xr}.
\label{eq:dcorr}
\end{equation}

\paragraph{(c) Re-projection (DFT).} Transform back with $\Gam^{+}$:
\begin{equation}
\boxed{\;
\frac{\partial\lamt}{\partial\xr}
=\Gam^{+}\,\bm m\,\Gam\,\bigl(\varepsilon\Ib-\Zr\bigr)\;}
\qquad\Longrightarrow\qquad
\frac{\partial\bm r}{\partial\xr}
=\Ib+\Yr\,\Gam^{+}\,\bm m\,\Gam\,\bigl(\varepsilon\Ib-\Zr\bigr).
\label{eq:Jfull}
\end{equation}

\paragraph{Why this is \emph{not} a holomorphic FRF map (the key
correctness point).} The factors $\Yr$ and $(\varepsilon\Ib-\Zr)$ are
complex-linear, but $\bm m$ is a real diagonal acting \emph{in the time
domain} between $\Gam$ and $\Gam^{+}$. The composite
$\Gam^{+}\bm m\,\Gam$ is therefore \emph{real-linear but not complex-linear}:
in the RI split it is a \emph{general} block
$\bigl[\begin{smallmatrix}\bm A_{RR}&\bm A_{RI}\\\bm A_{IR}&\bm A_{II}\end{smallmatrix}\bigr]$
with $\bm A_{RR}\neq\bm A_{II}$ and $\bm A_{RI}\neq-\bm A_{IR}$ in general.
Hence \eqref{eq:Jfull} \textbf{must} be assembled with the real operators
$\Gam,\Gam^{+}$ of \eqref{eq:dftentries}; collapsing it into the
holomorphic form $\bigl[\begin{smallmatrix}\Real&-\Imag\\\Imag&\Real\end{smallmatrix}\bigr]$
of the complex matrix $\Ib+\Yr(\varepsilon\Ib-\Zr)$ is only valid in
\emph{persistent} contact ($\bm m=\Ib$, where $\Gam^{+}\Gam=\Ib$) and is
wrong as soon as the contact set is partial. (This was the dominant bug in
the original \texttt{compute\_jacobian\_of\_residue\_RI}.) When $\bm m=\Ib$
the boxed result correctly collapses to
$\partial\bm r/\partial\xr=\varepsilon\,\Yr$.

\subsection{Derivative w.r.t.\ frequency}
Here $\xr$ is held fixed and the $\omega$-dependence enters through
$\Yb(\omega)$. Per-harmonic,
\begin{equation}
\frac{\partial\Zb_n}{\partial\omega}=-2n^2\omega\,\Mb+\mathrm i\,n\,\Cb,
\qquad
\frac{\partial\Yb_n}{\partial\omega}=-\Yb_n\,\frac{\partial\Zb_n}{\partial\omega}\,\Yb_n,
\label{eq:dY}
\end{equation}
and $\partial\Yr/\partial\omega=\Bb_{\mathrm c}(\partial\Yb/\partial\omega)\Bb_{\mathrm e}\T$,
$\partial\fadm/\partial\omega=\Bb_{\mathrm c}(\partial\Yb/\partial\omega)\tilde\fb_{\mathrm{ext}}$,
$\partial\Zr/\partial\omega=-\Zr(\partial\Yr/\partial\omega)\Zr$.
Differentiating \eqref{eq:residual},
\begin{equation}
\frac{\partial\bm r}{\partial\omega}
=\frac{\partial\Yr}{\partial\omega}\,\lamt
+\Yr\,\frac{\partial\lamt}{\partial\omega}
-\frac{\partial\fadm}{\partial\omega}.
\label{eq:drdw}
\end{equation}
The contact-force term uses the \emph{same} time-domain mask as
\eqref{eq:Jfull}; since $\partial\lamt_{\mathrm p}/\partial\omega$ is a
single (real-signal) vector its IDFT is unambiguous:
\begin{equation}
\frac{\partial\lamt}{\partial\omega}
=\Gam^{+}\,\bm m\,\Gam\;
\underbrace{\Bigl[\tfrac{\partial\Zr}{\partial\omega}(\fadm-\xr)+\Zr\,\tfrac{\partial\fadm}{\partial\omega}\Bigr]}_{\partial\lamt_{\mathrm p}/\partial\omega},
\quad
\frac{\partial\lamt_{\mathrm p}}{\partial\omega}
=\Zr\,\frac{\partial(\Bb\Yb)}{\partial\omega}\bigl(\tilde\fb_{\mathrm{ext}}-\Bb_{\mathrm e}\T\Zr(\fadm-\xr)\bigr).
\label{eq:dlamdw}
\end{equation}
The first two terms of \eqref{eq:drdw} combine to
$\Bb_{\mathrm c}(\partial\Yb/\partial\omega)(\Bb_{\mathrm e}\T\lamt-\tilde\fb_{\mathrm{ext}})$,
matching the implemented \texttt{compute\_derivative\_wrt\_omega\_RI}
(verified against finite differences to $\sim10^{-13}$).

% =====================================================================
\section{Properties, conditioning and algorithm}
% =====================================================================
\subsection{$\varepsilon$-independence of the solution}
At convergence ($\bm r=\bm0$) the solution is independent of
$\varepsilon$. With $\lambda=\lambda_{\mathrm p}+\lambda_{\mathrm c}$ and
$\lambda_{\mathrm c}=\max(0,-\lambda_{\mathrm p})\ge0$, substituting the
prediction \eqref{eq:pred} into \eqref{eq:residual} gives
$\bm r=\Yr(\lamt_{\mathrm c}-\varepsilon\tilde{\bm g})$, i.e.\ the impedance
residual $\Zr\bm r=\lamt_{\mathrm c}-\varepsilon\tilde{\bm g}$. The
correction satisfies $0\le\tfrac1\varepsilon\lambda_{\mathrm c}=g_{\mathrm p}\perp\lambda\ge0$
(Vadcard Eqs.~44--45), so the root $\tilde{\bm g}_{\mathrm p}=\tilde{\bm g}$
is reached independently of $\varepsilon$. Two practical caveats:
(i) $\varepsilon$ must be \emph{large enough} to detect contact through the
penalty -- in the test problem (stiffness $\sim1$) convergence requires
$\varepsilon\gtrsim10^{6}$ and is robust at $10^{8}$ (Vadcard sweeps
$\varepsilon\in[10^{8},2\!\cdot\!10^{12}]$); (ii) larger $\varepsilon$
amplifies the harmonic-truncation oscillations of the non-smooth
$\lambda(t)$, so the \emph{absolute} residual at the floor scales with
$\varepsilon$ -- the convergence test must be relative (below).

\subsection{Conditioning: why admittance vs.\ impedance matters}
\label{sec:cond-number}
The two residuals are related by $\bm r_{\mathrm{adm}}=\Yr\,\mathcal H_{\mathrm{imp}}$.
Near a substructure resonance $\Yr$ is large and ill-conditioned, so
admittance-form Newton is more sensitive there; this is intrinsic to FBS
with measured/numeric FRFs (van der Seijs after Eq.~3.26). The admittance
form is nevertheless required for experimental FRFs and is what the FBS
solver uses; the cure is a relative/scaled convergence criterion and arc-length
continuation with adaptive steps.

\subsection{Solver criterion}
Because of the truncation floor, the corrector accepts $\xr$ when
\emph{either} the residual is small relative to its value at the predictor,
$\|\bm r\|\le\eta_{\mathrm{rel}}\|\bm r^{(0)}\|$, \emph{or} the Newton step
stagnates, $\|\Delta\xr\|\le\eta_{\mathrm{stag}}\|\xr\|$ (the floor is
reached), in addition to the original absolute test
$\|\bm r\|<\eta_{\mathrm{abs}}$.

\subsection{Algorithm (drop-in replacement for AFT)}
\begin{center}
\fbox{\begin{minipage}{0.93\textwidth}
\textbf{DLFT residual/Jacobian at $(\xr,\omega)$} -- same interface as AFT:
\begin{enumerate}[leftmargin=1.5em,itemsep=1pt]
\item Build $\Yb(\omega)$ (numeric from $\Mb,\Cb,\Kb$, or interpolated FRF);
$\Yr=\Bb_{\mathrm c}\Yb\Bb_{\mathrm e}\T$, $\fadm=\Bb_{\mathrm c}\Yb\tilde\fb_{\mathrm{ext}}$, $\Zr=\Yr\inv$.
\item \emph{Prediction:} $\lamt_{\mathrm p}=\Zr(\fadm-\xr)+\varepsilon(\xr-\tilde{\bm g}_0)$;
$\lambda_{\mathrm p}=\Gam\lamt_{\mathrm p}$ \hfill (IDFT).
\item \emph{Correction:} $\lambda=\max(0,\lambda_{\mathrm p})$, mask $\bm m$;
$\lamt=\Gam^{+}\lambda$ \hfill (DFT).
\item \emph{Residual:} $\bm r=\xr+\Yr\lamt-\fadm$.
\item \emph{Jacobian:}
$\partial\bm r/\partial\xr=\Ib+\Yr\,\Gam^{+}\bm m\,\Gam(\varepsilon\Ib-\Zr)$;
$\partial\bm r/\partial\omega$ from \eqref{eq:drdw}--\eqref{eq:dlamdw}.
\end{enumerate}
The corrector solves $\bm r(\xr)=\bm0$ for $\xr$ inside the
predictor--corrector arc-length continuation in $\omega$; \emph{only}
steps 2--3 (AFT $\fb_{\mathrm{nl}}$) and step 5 differ between AFT and DLFT.
Response recovery uses
$\hat\ub=\Yb(\tilde\fb_{\mathrm{ext}}-\Bb_{\mathrm e}\T\lamt)$
(Soleimani Eqs.~25--26).
\end{minipage}}
\end{center}

% =====================================================================
\section{Map to the \texttt{pyhbm} implementation}
% =====================================================================
\begin{center}\small
\begin{tabular}{@{}lll@{}}
\toprule
Symbol & Meaning & \texttt{FBS\_DLFT} member \\
\midrule
$\xr$ & rel.\ displacement coeff.\ (Newton unknown) & \texttt{x.fourier.coefficients} \\
$\Bb_{\mathrm c}=\Bb_{\mathrm e}=\Bb$ & signed Boolean (multiharmonic) & \texttt{self.B\_fourier} \\
$\Yb$ & block-diag substructure admittance & \texttt{compute\_FRF(omega)} \\
$\Yr=\Bb\Yb\Bb\T$ & interface admittance & \texttt{\_get\_Yr} \\
$\fadm=\Bb\Yb\tilde\fb_{\mathrm{ext}}$ & free relative response & \texttt{\_get\_Fext\_admr} \\
$\Zr(\fadm-\xr)$ & balancing force (pre-penalty) & \texttt{\_get\_Zr\_rhs} \\
$\lamt=\Gam^{+}\max(0,\lambda_{\mathrm p})$ & corrected contact force & \texttt{\_get\_lambda\_corrected} \\
$\bm m$ & contact mask $\mathbf1[\lambda_{\mathrm p}>0]$ & \texttt{x.contact\_mask} \\
$\Gam,\Gam^{+}$ & real IDFT/DFT \eqref{eq:dftentries} & \texttt{self.\_idft\_RI},\,\texttt{self.\_dft\_RI} \\
$\partial\bm r/\partial\xr$ & Eq.~\eqref{eq:Jfull} & \texttt{compute\_jacobian\_of\_residue\_RI} \\
$\partial\bm r/\partial\omega$ & Eq.~\eqref{eq:drdw} & \texttt{compute\_derivative\_wrt\_omega\_RI} \\
\bottomrule
\end{tabular}
\end{center}

\paragraph{Diagnosed defects of the original code (corrected per this note).}
(b) the matrix $\partial\lamt/\partial\xr$ was built with \texttt{rfft/irfft}
on \emph{complex} columns and then wrapped by the holomorphic
\texttt{FRF\_to\_RI}, i.e.\ \eqref{eq:Jfull} assembled as if
$\Gam^{+}\bm m\,\Gam$ were complex-linear -- wrong for partial contact;
(c) $\varepsilon=10^{2}$ was orders of magnitude too small;
(a) the sign of the penalty is self-consistent with a wall at $+g_0$
(equivalent to Vadcard up to the sign of $x_{\mathrm r}$);
(d) the time-domain correction \eqref{eq:corr} before the DFT was already
correct; $\partial\bm r/\partial\omega$ was already correct.

\begin{thebibliography}{9}
\bibitem{Vadcard2022} T.\ Vadcard, A.\ Batailly, F.\ Thouverez,
\emph{On Harmonic Balance Method-based Lagrangian contact formulations for
vibro-impact problems}, J.\ Sound Vib.\ \textbf{531} (2022) 116950.
\bibitem{VdS2016} M.\ van der Seijs, \emph{Experimental Dynamic
Substructuring}, PhD thesis, TU Delft (2016), Ch.~3 (LM-FBS, Eqs.~3.20--3.28).
\bibitem{Soleimani2025} H.\ Soleimani, N.\ Aage,
\emph{A closed-form nonlinear dynamic substructuring method}, CMAME
\textbf{439} (2025) 117882 (Eq.~24).
\bibitem{Krack2019} M.\ Krack, J.\ Gross, \emph{Harmonic Balance for
Nonlinear Vibration Problems}, Springer (2019).
\end{thebibliography}

\end{document}
